\section{Inspiration in der Bibel}
Nachdem wir uns am letzten Sonntag mit den verschiedenen Theorien der Inspiration beschäftigt haben, wollen wir Heute anschauen was denn überhaupt Inspiration ist und was die Schrift dazu sagt. Wenn wir über die Inspiration reden wollen, kommen wir nicht an der Aussage von Paulus im 2. Timotheus Brief vorbei. Lesen wir zum Anfang diesen Vers.
\begin{bibelbox}{SCHL}{IITim}{3:16}
...
\end{bibelbox}
Aber nicht nur dieser Vers zeigt, dass Gott der Urheber der Schrift ist, sondern auch Petrus in seinem zweiten Brief ermutigt die Empfänger des Briefes sich an der Schrift zu orientieren und sagt, dass diese nicht von Menschen geschrieben wurde, sondern durch den Heiligen Geist sprachen die Menschen.
\begin{bibelbox}{SCHL}{IIPetr}{1:21}
...
\end{bibelbox}  
Inspiration kommt aus dem lateinischen \textit{inspirata}. In der Vulgata wurde die Aussage \enquote{von Gott gehaucht} mit inspirata übersetzt. Darum wird dieses Wort heute benutzt um dieses \enquote{von Gott gehaucht} darzustellen.

Es gibt viele Stellen in der Bibel die uns aufzeigen, dass die Bibel Gottes Wort ist. Im Matthäus Evangelium gibt es viele Zitate die auf das Alte Testament verweisen. Auch Jesus bezeugt in 
\begin{bibelbox}{SCHL}{Matt}{5:18}
...
\end{bibelbox}
Ich denke das ist \textit{\enquote{titsch und ditlich}} wie wir hier im Wallis sagen. Aber wie das Alte Testament als von Gott eingehauchte Schrift dargestellt wird, so auch das Neue Testament. Petrus zitiert im Vers \biblezit{ITim}{5:18} wird im gleichen Atemzug aus dem alten Testament \biblezit{VMos}{25:4} und aus dem neuen Testament \biblezit{Lk}{10:7} zitiert.
\begin{bibelbox}{SCHL}{ITim}{5:18}
...
\end{bibelbox}
Aber auch Jesus sagt, dass die Schreiber vom Neuen Testament vom Heiligen Geist geführt sind.
\begin{bibelbox}{SCHL}{Joh}{16:12-13}
...
\end{bibelbox}

Und auch Petrus schreibt in
\begin{bibelbox}{SCHL}{IIPetr}{3:15-16}
...
\end{bibelbox}
das die paulinischen Briefe auch zur Schrift gehöhren. Also auch zum inspirierten Wort.

Die Erkentniss, dass die Schrift von Gott gehaucht ist, ist nicht neu. Schon Augustinus schreibt:
\begin{quote}
    Jene Schrift, die nicht durch zufälligen Eindruck auf die Gemüter, sondern geradezu kraft Anordung der höchstenn Vorhersehung mehr als alles übrige SChriftum welchen Volkes auch immer die Menschengeister jeder Richtung durch die sie auszeichnende göttliche Autorität in ihrem Banne hält.
\end{quote}

Aber auch die katholische Kirche sagt, dass sich die Schrift sich selbst bezeugt, oder \enquote{selbst-treu} ist -- Theologisch \textit{Autopistos}.
So wird in einem wichtigen Text der katholischen Kirche fogendes geschrieben:
\begin{quote}
    Zur Abfassung der Heiligen Bücher hat Gott Menschen erwählt, die ihm durch den Gebrauch ihrer eigenen Fähigkeiten und Kräfte dazu dienen sollten, all das und nur das, was er -- in ihnen und durch sie wirksam -- geschrieben haben wollte, als echte Verfasser schriftlich zu überliefern
\end{quote}

Aber egal was Theologen, Kirchen und Kirchenväter sagen, wir halten uns an den Aussagen in der Schrift selber, darum nochmals 
\begin{bibelbox}{SCHL}{IITim}{3:16}
...
\end{bibelbox}